\section{Introduction}
\label{sec:intro}

Aspirin (\ce{C9H8O4}) is synthesised by esterification of the phenolic hydroxyl
group of salicylic acid with acetic anhydride, using a few drops of concentrated
sulfuric acid as a catalyst~\cite{clayden2012,pavia2015}. The overall reaction is
shown in Equation~\eqref{eq:overall}:
\begin{equation}
  \ce{C7H6O3 + (CH3CO)2O ->[\ce{H2SO4}] C9H8O4 + CH3COOH}
  \label{eq:overall}
\end{equation}

The structure of the product is shown below.

\begin{center}
  \chemfig{*6((-COOH)=-=-=(-O-[::60]C(=[::60]O)-[::-60]CH_3)-)}
\end{center}

Excess acetic anhydride is destroyed during work-up by hydrolysis:
\ce{(CH3CO)2O + H2O -> 2 CH3COOH}. Unreacted salicylic acid gives a violet
complex with iron(III), \ce{Fe^3+ + 3 C7H5O3- -> [Fe(C7H5O3)3]}, which is the
basis of the purity test described in Section~\ref{sec:procedure}.
